\section{Общая архитектура программного комплекса на языке Rust}

Программный комплекс реализован по принципу
«одно вычислительное ядро и несколько каналов доставки результата».
Все термодинамические вычисления сосредоточены в \texttt{if97\_core}.
Пользовательские приложения и сервисный контур
используют единый программный интерфейс ядра и не дублируют формулы IF97.

Такое разделение решает две инженерные задачи.
Первая задача --- совпадение результата во всех каналах доставки.
Вторая задача --- локализация математической логики в одном модуле,
что упрощает сопровождение и верификацию.

Для обмена данными применяется общий слой DTO \texttt{if97\_app\_api}.
Он фиксирует контракт входов и выходов между компонентами.
Реализация на Rust обеспечивает строгую типизацию,
контролируемую обработку ошибок и высокую производительность
без потери надёжности~\cite{rust_book}.

Математическая основа вычислительного ядра опирается на два фундаментальных потенциала.
В регионах 1, 2 и 5 используется безразмерная энергия Гиббса
\begin{equation}
\gamma(\pi,\tau)=\frac{g(p,T)}{RT},
\end{equation}
а в регионе 3 --- безразмерная энергия Гельмгольца
\begin{equation}
\phi(\delta,\tau)=\frac{f(\rho,T)}{RT}.
\end{equation}
Через производные \(\gamma_{\pi}\), \(\gamma_{\tau}\), \(\gamma_{\tau\tau}\)
и \(\phi_{\delta}\), \(\phi_{\tau}\), \(\phi_{\delta\delta}\)
ядро вычисляет удельный объём, энтальпию, энтропию, теплоёмкость и скорость звука,
поэтому выбор региона непосредственно определяет набор используемых формул.

Чтобы показать фактический путь данных через внешний API, валидацию и маршрутизацию ядра,
ниже приведена схема вычислительного контура.

\begin{figure}[htbp]
\centering
\begin{tikzpicture}[
    >=Stealth,
    font=\scriptsize,
    block/.style={draw, rounded corners, align=center, minimum height=8mm, text width=3.1cm, fill=blue!6},
    api/.style={block, fill=green!9},
    core/.style={block, fill=orange!12},
    route/.style={draw, rounded corners, align=center, minimum height=8mm, text width=3.1cm, fill=gray!8},
    branch/.style={draw, rounded corners, align=center, minimum height=8mm, text width=2.4cm, fill=gray!5},
    conv/.style={draw, rounded corners, align=center, minimum height=6mm, text width=1.8cm, fill=yellow!10},
    region/.style={draw, rounded corners, align=center, minimum height=10mm, text width=2.4cm, fill=blue!4},
    result/.style={draw, rounded corners, align=center, minimum height=10mm, text width=4.0cm, fill=blue!8},
    error/.style={draw, rounded corners, align=center, minimum height=18mm, text width=3.1cm, fill=red!12},
    % Новые цвета маршрутов
    pt_style/.style={draw, thick, dashed, orange!90!black},
    phps_style/.style={draw, thick, dashed, violet!80!black},
    rhot_style/.style={draw, thick, dashed, green!70!black},
    px_style/.style={draw, thick, dashed, blue!80!black}
]
    \node[block] (client) at (0,0) {Внешний\\потребитель\\FLTK / Tauri / gRPC};
    \node[api] (dto) at (0,-1.35) {if97\_app\_api\\DTO и схема входа};
    \node[core] (engine) at (0,-2.8) {if97\_core::engine\\валидация и регион};
    \node[route] (split) at (0,-4.15) {Выбор маршрута};

    % Блок ошибок теперь расположен компактнее
    \node[error] (err) at (5.0,-2.075) {Ошибка входа, схемы,\\расчёта или региона};

    % Уменьшенный шаг между колонками для большей компактности
    \node[branch] (pt) at (-4.35,-5.6) {$pT$\\вход};
    \node[branch] (phps) at (-1.45,-5.6) {$ph/ps$\\вход};
    \node[branch] (rhot) at (1.45,-5.6) {$\rho T$\\вход};
    \node[branch] (px) at (4.35,-5.6) {$px$\\вход};

    \node[conv] (sec_p1) at (-5.8,-7.6) {$p$ секущей};
    \node[conv] (sec_p2) at (-2.9,-7.6) {$p$ секущей};
    \node[conv] (sec_rho3) at (0,-7.6) {$\rho$ секущей};
    \node[conv] (calc_x4) at (2.9,-7.6) {вычисл. $x$};
    \node[conv] (sec_p5) at (5.8,-7.6) {$p$ секущей};

    \node[region] (r1) at (-5.8,-9.2) {Регион 1\\$pT, ph, ps$};
    \node[region] (r2) at (-2.9,-9.2) {Регион 2\\$pT, ph, ps$};
    \node[region] (r3) at (0,-9.2) {Регион 3\\$\rho T, ph, ps$};
    \node[region] (r4) at (2.9,-9.2) {Регион 4\\$px$};
    \node[region] (r5) at (5.8,-9.2) {Регион 5\\$pT, ph, ps$};

    \node[result] (state) at (0,-11.2) {if97\_core::domain::WaterState\\полный набор свойств};
    \node[api] (ser) at (0,-12.5) {if97\_app\_api\\результат или ошибка};
    \node[block] (resp) at (0,-13.8) {Ответ клиенту};

    \draw[->] (client) -- (dto);
    \draw[->] (dto) -- (engine);
    \draw[->] (engine) -- (split);

    \draw[->] (split) -- (pt);
    \draw[->] (split) -- (phps);
    \draw[->] (split) -- (rhot);
    \draw[->] (split) -- (px);

    \coordinate (bus_pt) at (-4.35, -6.2);
    \coordinate (bus_phps) at (-1.45, -6.4);
    \coordinate (bus_rhot) at (1.45, -6.6);

    % pT - оранжевый
    \draw[pt_style, -] (pt.south) -- (bus_pt);
    \draw[pt_style, ->] (bus_pt) -| ([xshift=-10pt]r1.north);
    \draw[pt_style, ->] (bus_pt) -| ([xshift=-10pt]r2.north);
    \draw[pt_style, ->] (bus_pt) -| (sec_rho3.north);
    \draw[pt_style, ->] (bus_pt) -| ([xshift=-10pt]r5.north);

    % Ошибка pT: проходит аккуратно слева и между DTO и Engine
    \draw[->, red!70, dashed, thick] (pt.west) -- ++(-2.0,0) |- (err.west); 

    % phps - фиолетовый
    \draw[phps_style, -] (phps.south) -- (bus_phps);
    \draw[phps_style, ->] (bus_phps) -| ([xshift=10pt]r1.north);
    \draw[phps_style, ->] (bus_phps) -| ([xshift=10pt]r2.north);
    \draw[phps_style, ->] (bus_phps) -| ([xshift=-10pt]r3.north);
    \draw[phps_style, ->] (bus_phps) -| ([xshift=10pt]r5.north);

    % rhot - зелёный
    \draw[rhot_style, -] (rhot.south) -- (bus_rhot);
    \draw[rhot_style, ->] (bus_rhot) -| (sec_p1.north);
    \draw[rhot_style, ->] (bus_rhot) -| (sec_p2.north);
    \draw[rhot_style, ->] (bus_rhot) -| ([xshift=10pt]r3.north);
    \draw[rhot_style, ->] (bus_rhot) -| (calc_x4.north);
    \draw[rhot_style, ->] (bus_rhot) -| (sec_p5.north);

    % px - синий
    \draw[px_style, ->] (px.south) -- ++(0,-1.0) -| ([xshift=10pt]r4.north);

    % Линии от конвертеров к регионам - сплошные, тонкие (без стиля, чтобы не перегружать)
    \draw[->] (sec_p1.south) -- (r1.north);
    \draw[->] (sec_p2.south) -- (r2.north);
    \draw[->] (sec_rho3.south) -- (r3.north);
    \draw[->] (calc_x4.south) -- (r4.north);
    \draw[->] (sec_p5.south) -- (r5.north);

    % Сборка результата
    \draw[->] (r1.south) |- ([yshift=15pt]state.north) -- (state.north);
    \draw[->] (r2.south) |- ([yshift=15pt]state.north) -- (state.north);
    \draw[->] (r3.south) -- (state.north);
    \draw[->] (r4.south) |- ([yshift=15pt]state.north) -- (state.north);
    \draw[->] (r5.south) |- ([yshift=15pt]state.north) -- (state.north);

    \draw[->] (state) -- (ser);
    \draw[->] (ser) -- (resp);

    % Линии ошибок - красные пунктирные
    \draw[->, red!70, dashed, thick] (dto.east) -- (err.west |- dto.east);
    \draw[->, red!70, dashed, thick] (engine.east) -- (err.west |- engine.east);
    % Эта линия спускается строго по правой границе, никого не задевая
    \draw[->, red!70, dashed, thick] (err.east) -| (7.3, -12.5) -- (ser.east);

\end{tikzpicture}
\caption{Путь данных через слой DTO, валидацию и вычислительное ядро}
\label{fig:core-data-path}
\end{figure}
